\begin{figure}[!htb] 
	\centering
	\resizebox{\textwidth}{!}{
		\subcaptionbox{Trajectoire en force de la première session i) \label{subfig_force_filtering_a}}
		{
			\tikzsetnextfilename{fig_force_filtering_a}
			\begin{tikzpicture}
			\centering
			%
			\begin{axis}[%
				width=.70\linewidth,
				height=.25\linewidth,
				scale only axis,
				xlabel={temps (\SI{}{\second})},
				ylabel={Force (\SI{}{\newton})},
				axis x line = bottom,
				axis y line = left, 
				]
				%
				\addplot [
					color=blue,
					solid,
				]
				table [col sep=comma, x=time, y=force] {misc/force_signal_filtered_chunk.csv};
				%
				\addplot [
					color=red,
					smooth,
					]
				table [col sep=comma, x=time, y=filtered] {misc/force_signal_filtered_chunk.csv};
				\addlegendentry{mesure};
				\addlegendentry{filtré};
				%
			\end{axis}		
			\end{tikzpicture}
		}
		\subcaptionbox{Trajectoire en force avec retours haptiques et perturbations aléatoires\label{subfig_force_filtering_b}}
		{
			\tikzsetnextfilename{fig_force_filtering_b}
			\begin{tikzpicture}
				\centering
				%
				\begin{axis}[%
					width=.25\linewidth,
					height=.25\linewidth,
					scale only axis,
					xlabel={temps (\SI{}{\second})},
					axis x line = bottom,
					axis y line = left, 
					]
					%
					\addplot [
						color=blue,
						solid,
						forget plot,
						]
					table [col sep=comma, x=time, y=force] {misc/force_signal_perturbed_filtered_chunk.csv};
					%
					\addplot [
						color=red,
						smooth,
						forget plot,
						]
					table [col sep=comma, x=time, y=filtered] {misc/force_signal_perturbed_filtered_chunk.csv};
					%
					\addplot [
						only marks,
						color=black,
						mark=triangle,
						]
					table [col sep=comma, x=time, y=perturbation] {misc/force_signal_perturbed_filtered_chunk.csv};
					\addlegendentry{pert.};
				%
				\end{axis}		
			\end{tikzpicture}
		}
	}	
	\caption{Exemples du filtrage passe-bas décrit en \cref{chap3_sec_method_sub_virt_traj_4}, pour des trajectoires en force.}
	\label{fig_force_filtering}
\end{figure}